%% HAND-WRITTEN fragment -- NOT generated by maestro2tex (XY figure kind).
%%
%% Data: data/electrode_panel_{ua,aa,h2o2,pyo}.dat, written by
%% ~/chips/castalia/skill/cv_panel_fig.py --dat from the Maestro corner run
%% archived at ~/chips/castalia/simruns/cv_panel (tb_anatop_pixel/CvBvTB, test
%% ps_cv, analyte corners UricAcid/Ascorbate/H2O2_PB/Pyocyanin defined in that
%% maestro view; 5/5 completed, summary_ps_cv.txt). Columns: E = V(RE)-V_CM [V],
%% i_we = (V_CM - V_OUT)/R_f [uA] (anodic positive). Same data as ISCAS27
%% Fig. 4. Re-run the corner set and the script to regenerate.
\providecommand{\MaestroRoot}{include/analog/}
\providecolor{mtxA}{HTML}{2A78D6}
\providecolor{mtxB}{HTML}{EB6834}
\providecolor{mtxC}{HTML}{1BAF7A}
\providecolor{mtxD}{HTML}{EDA100}
\begin{figure}[htbp]
  \centering
  {\pgfplotsset{compat=1.18}%
  \begin{tikzpicture}
    \begin{axis}[
      width=0.78\linewidth, height=6.4cm,
      grid=both,
      major grid style={line width=0.2pt, draw=black!14},
      minor grid style={line width=0.2pt, draw=black!7},
      minor tick num=1,
      axis line style={line width=0.4pt, draw=black!45},
      tick style={line width=0.4pt, draw=black!45},
      tick align=outside,
      label style={font=\small}, tick label style={font=\small},
      legend style={font=\footnotesize, draw=none, fill=white,
                    fill opacity=0.85, text opacity=1,
                    at={(0.03,0.97)}, anchor=north west, row sep=1pt},
      every axis plot/.append style={line join=round},
      xlabel={\(E = V_{\mathrm{RE}} - V_{\mathrm{CM}}\)~[\si{\volt}]},
      ylabel={\(i_{\mathrm{we}}\)~[\si{\micro\ampere}]},
      enlarge x limits=0.03, enlarge y limits=0.08,
    ]
      \addplot[mtxA, line width=1.0pt, no marks]
        table {\MaestroRoot data/electrode_panel_ua.dat};
      \addlegendentry{uric acid, \SI{500}{\micro\mole\per\litre}}
      \addplot[mtxB, line width=1.0pt, no marks]
        table {\MaestroRoot data/electrode_panel_aa.dat};
      \addlegendentry{ascorbate, \SI{100}{\micro\mole\per\litre}}
      \addplot[mtxC, line width=1.0pt, no marks]
        table {\MaestroRoot data/electrode_panel_h2o2.dat};
      \addlegendentry{H$_2$O$_2$/PB, \SI{60}{\micro\mole\per\litre}}
      \addplot[mtxD, line width=1.0pt, no marks]
        table {\MaestroRoot data/electrode_panel_pyo.dat};
      \addlegendentry{pyocyanin, \SI{50}{\micro\mole\per\litre}}
      \addplot[black!45, line width=0.4pt, domain=-0.42:0.82, samples=2,
        forget plot] {0};
    \end{axis}
  \end{tikzpicture}}
  \caption[{Four-analyte voltammograms through the rev-2 channel.}]{%
  Voltammograms of the four-analyte wound panel through the rev-2 channel, one
  electrode parameter set per channel, from the cyclic-voltammetry
  Maestro corner run of Section~\ref{ss:analog-electrode-bc}
  (\SI{2}{\volt\per\second}, \(R_f = \SI{35}{\kilo\ohm}\), conditions of the
  in-loop run above). Peak potential identifies the analyte and peak height its
  concentration; each trace is one full cycle started at that analyte's
  non-faradaic vertex. \(E^{0\prime}\) and concentrations are
  literature-anchored; \(k^{0}\) is fitted (no measured value exists on carbon
  for any of the four), and reversal peaks appear for the chemically
  irreversible couples because following chemistry is not modelled. The two
  cathodic markers span \SIrange{-1.4}{-2.4}{\micro\ampere} --- unmeasurable by
  the unidirectional rev-1 front end.}
  \label{fig:electrode-panel}
\end{figure}
